\section{Flattening with the New, More Efficient Rules}
My code validates on default input (which is checking the implementation
against the classic version) and the implementation is provided below:
\begin{lstlisting}[language={futhark}]
def optimII1Ker [m][n][q] (Sa: [m]u32, Da: [n]f32)
                          (Sb: [m]u32, Db: [q]f32)
                          (cs: [m]f32) (inds: [m]i64)
                        : [m]f32 =
  let (Ba, flags'') = mkFlagArray Sa 0 (iota m) 
  let flags' = map u32.i64 flags''
  let flags = map bool.u32 flags' :> [n]bool
  let flen_iot = reduce (+) 0 Sa |> i64.u32 |> iota :> [n]i64
  let iia1 = sgmScan (+) 0 flags (flags' :> [n]u32)
  let iia2 = map2(\i sgm -> (u32.i64 i) - Ba[i64.u32 sgm]) (flen_iot) iia1 
  let (Bb, _) = mkFlagArray Sb false (replicate m true)
  let b_vals = map2 (\off ind -> Db[i64.u32 off + ind]) Bb inds
  let tmp1s = map2 (\a sgm -> (f32.sqrt a) * b_vals[i64.u32 sgm] + cs[i64.u32 sgm]) Da iia1 
  let iotis = iia2 
  let iotfs = map f32.u32 iotis
  let tmp2s = map2 (+) tmp1s iotfs
  let tmp_scan = sgmScan f32.max f32.lowest flags tmp2s
  let res = imap2 (iota m) Sa 
          (\i s -> if s <= 0 then f32.lowest 
                   else if i == m-1
                        then tmp_scan[n-1]
                        else tmp_scan[i64.u32 Ba[i+1]-1]
          ) 
  in  res
\end{lstlisting}

The runtime comparison between the standard classic implementation and the new
optimized implementation can be seen below:
\begin{table}[ht]
\centering
\begin{tabular}{c|c|c|c}
\hline
\textbf{(m, q, p)} & \textbf{Classic ($\mu s$)} & \textbf{optimII1 ($\mu s$)} & \textbf{Speedup} \\
\hline
(10, 10\,000\,000, 100\,000)     & 11212 & 7808 & 1.44 \\
(100, 1\,000\,000, 10\,000)      & 11192 & 7810 & 1.43 \\
(1000, 100\,000, 1000)           & 11250 & 7807 & 1.44 \\
(10000, 10\,000, 100)            & 11269 & 7866 & 1.45 \\
(100000, 1000, 10)               & 11340 & 7908 & 1.43 \\
\hline
\end{tabular}
\caption{Performance comparison of classic and optimII1 kernels with speedup.}
\end{table}

This means that the new optimized implementation improves performance by around
1.44 across the board. This follows our general intuition about how the new
flattening rules reduce the overall memory overhead needed and thus leads to
performance improvements.

\subsection*{Note on running the tests}
The built in tests require the datasets to be generated first via. a call to
\texttt{make datasets} within the subfolder. Then one can call \texttt{make
run} to run the program and seeing that it validates.

